\documentclass[11pt]{amsart}
\usepackage{amsmath,amssymb,amsthm}
\usepackage{booktabs}
\usepackage{hyperref}
\usepackage{microtype}
\emergencystretch=1em

\newtheorem{theorem}{Theorem}
\newtheorem{corollary}[theorem]{Corollary}
\newtheorem{proposition}[theorem]{Proposition}
\theoremstyle{definition}
\newtheorem{definition}[theorem]{Definition}
\theoremstyle{remark}
\newtheorem{remark}[theorem]{Remark}

\title[Rational Contamination in PCF Searches]{Trivial rational contamination in PSLQ-based
  polynomial continued fraction searches:\\
  diagnosis, correction, and a pre-screening protocol}

\author{Papanokechi}
\address{Independent researcher}
\email{papanokechi@proton.me}

\subjclass[2020]{11A55, 11Y65, 11B37}
\keywords{polynomial continued fractions, PSLQ, false positives,
  rational contamination, Wallis recurrence}

\date{\today}

\thanks{This research used AI assistance (Claude, Anthropic) for
  code generation, conjecture exploration, and manuscript drafting
  within an AEAL (Adversarial Epistemic Audit Loop) relay framework.
  All mathematical claims were verified by the author independently.}

\begin{document}
\begin{abstract}
We identify two mechanisms that produce rational limits in random
polynomial continued fractions (PCFs), both capable of generating
false positives in PSLQ-based searches for transcendental identities.
First, if $a(1)=0$ then the GCF $b(0)+\mathbf{K}(a(n)/b(n))$
reduces trivially to $b(0)$ (the first partial numerator vanishes).
Second, if $a(k)=0$ for some $k\ge 2$, the CF terminates at finite
depth.
Together these account for \emph{all} 43 rational limits in a
1000-family survey across degrees~2--6.
We propose a pre-screening protocol that eliminates both classes
before extended-basis PSLQ, and document the five-iteration
discovery arc---including a Wallis initialization bug that temporarily
masked the triviality of the first mechanism---as a case study
in machine-assisted self-correction.
\end{abstract}

\maketitle

% =====================================================================
\section{Introduction}\label{sec:intro}
% =====================================================================

A \emph{polynomial continued fraction} (PCF) is a generalized continued
fraction
\begin{equation}\label{eq:pcf}
  K \;=\; b(0) + \cfrac{a(1)}{b(1) + \cfrac{a(2)}{b(2) + \cfrac{a(3)}{\ddots}}}
\end{equation}
where $a(n) = \sum_{i=0}^{d} \alpha_i\, n^i$ and
$b(n) = \sum_{i=0}^{d} \beta_i\, n^i$
are polynomials of degree~$d$ with integer coefficients.
Following the Ramanujan Machine project~\cite{raayoni2021}, random sampling of
such families has become a standard tool for discovering new identities
involving mathematical constants.

Any PSLQ-based search for transcendental identities among PCF limits
must contend with a practical nuisance: a non-negligible fraction of
random families converge to \emph{rational} values, which can trigger
PSLQ relations that mimic genuine transcendental identities when the
basis includes redundant constants.
In a survey of 1000 random PCF families (200 per degree $d=2,\ldots,6$,
coefficients in $\{-4,\ldots,4\}$), we found 43 such families---initially
misidentified as evidence for $\ln 2$ universality.

The correct explanation is elementary:

\begin{theorem}[Trivial rational limit]\label{thm:main}
  Let $a(n),b(n)\in\mathbb{Z}[n]$.
  If $a(1)=0$, then $K = b(0)$.
\end{theorem}

\begin{proof}
  The first partial numerator $a(1)$ is the numerator of the entire tail
  $a(1)/(b(1)+\cdots)$.  When $a(1)=0$, the tail vanishes and
  $K = b(0) + 0 = b(0)$.
\end{proof}

The condition $a(1)=0$ means $\sum_{i=0}^d \alpha_i = 0$---the sum
of the coefficients of the numerator polynomial vanishes.
This is a linear constraint on the $d+1$ coefficients and is satisfied
by $\sim\!5$--$9\%$ of random families depending on degree
(Section~\ref{sec:density}).
Although elementary, this criterion identifies a non-negligible
fraction of random PCF families as trivially rational---sufficient
to dominate any PSLQ search that does not filter them and to produce
false-positive transcendental identity claims.
We state it as a theorem rather than a remark precisely because its
role is \emph{diagnostic}: it is the first step in the pre-screening
protocol of Section~\ref{sec:protocol}.

\begin{proposition}[Finite CF]\label{prop:finite}
  If $a(k)=0$ for some integer $k\ge 2$ and $a(n)\ne 0$ for $1\le n<k$,
  then the GCF\,\eqref{eq:pcf} terminates at depth $k-1$ and is
  rational.
\end{proposition}

\begin{proof}
  The $k$-th partial numerator vanishes, so the tail from depth~$k$
  onward contributes zero.  The result is a depth-$(k-1)$ finite
  continued fraction, which is a rational function of its entries.
\end{proof}

Together, Theorem~\ref{thm:main} and Proposition~\ref{prop:finite}
explain all 43 rational limits in our survey (Table~\ref{tab:classification}),
with zero residual cases.

% =====================================================================
\section{The Wallis recurrence and a convention warning}\label{sec:wallis}
% =====================================================================

The standard Wallis recurrence for the convergents of~\eqref{eq:pcf} is
\begin{equation}\label{eq:wallis}
  A_{-1}=1,\quad A_0 = b(0),\qquad
  B_{-1}=0,\quad B_0 = 1,
\end{equation}
with
\begin{equation}\label{eq:wallis-rec}
  A_n = b(n)\,A_{n-1} + a(n)\,A_{n-2},\qquad
  B_n = b(n)\,B_{n-1} + a(n)\,B_{n-2},
\end{equation}
for $n\ge 1$.  The $n$-th convergent is $K_n = A_n/B_n$.

An alternative implementation computes the
\emph{tail} CF $T = a(1)/(b(1)+\cdots)$ separately and returns
$K = b(0) + T$.  For the tail, the correct Wallis initialization is
$P_0 = a(1),\, Q_0 = b(1)$, giving $T = P_n/Q_n$.

\begin{remark}[The initialization bug]\label{rem:bug}
  Our original evaluator used the incorrect initialization
  \[
    \texttt{Pc} = b(1),\quad \texttt{Qc} = 1
    \qquad\text{(wrong)}
  \]
  instead of the correct
  \[
    \texttt{Pc} = b(0),\quad \texttt{Qc} = 1,\quad
    \text{with } K_n = A_n/B_n
    \qquad\text{(standard Wallis)}
  \]
  and returned $b(0) + \texttt{Pc}/\texttt{Qc}$.
  This computes $b(0) + [b(1) + T]$ instead of $b(0) + T$, adding
  an offset of exactly $b(1)$ to every PCF value.

  \smallskip\noindent\textbf{Concrete example.}
  Let $a(n) = n$ and $b(n) = 2+n$ (so $a(1)=1$, $b(0)=2$, $b(1)=3$).
  The CF is $2 + 1/(3 + 2/(4 + 3/(5+\cdots)))$.
  The standard Wallis gives $K \approx 2.2956$;
  the bugged evaluator returns $2.2956 + 3 = 5.2956$.

  \smallskip
  The error is invisible to convergence testing
  (the CF still converges, just to the wrong value) and is not
  caught by PSLQ rationality checks (a shifted rational is still
  rational).  It was detected only when the ``theorem''
  $K = b(0)+b(1)$ was subjected to a convention audit against
  the golden ratio $\varphi = 1+1/(1+1/(1+\cdots))$
  (Section~\ref{sec:discovery}, Iteration~11).
  The bug has been patched; the corrected evaluator and the
  patch commit are recorded in the project repository.%
  \footnote{Bug patch: \texttt{siarc\_t1t6\_relay.py}, function
    \texttt{eval\_pcf\_forward}. All results in this paper use
    the corrected standard Wallis~\eqref{eq:wallis}--\eqref{eq:wallis-rec}.}
\end{remark}

We emphasize this because the bug is easy to introduce and hard
to detect: any code that splits the CF into $b(0)$ plus a tail
must take care with the tail's initial conditions.
The standard formulation~\eqref{eq:wallis}--\eqref{eq:wallis-rec}
avoids the issue entirely.

% =====================================================================
\section{Empirical classification}\label{sec:empirical}
% =====================================================================

\subsection{Survey design}
For each degree $d\in\{2,3,4,5,6\}$, we generated 200 random PCF families
with coefficients $\alpha_i,\beta_i$ drawn uniformly from $\{-4,\ldots,4\}$
(leading coefficients and $\beta_0$ forced nonzero).
Each family was evaluated using the standard Wallis
recurrence~\eqref{eq:wallis}--\eqref{eq:wallis-rec} to depth~2000
at 150-digit precision, and the limit~$K$ was tested for rationality
via PSLQ on the basis $[1,K]$.%
\footnote{Reproduce:
  \texttt{python scripts/iter11\_bugfix\_reeval.py}.
  Full code at \url{https://github.com/papanokechi/pcf-rational-contamination}.}

\subsection{Results}

\begin{table}[ht]
\centering
\caption{Classification of the 43 rational PCF limits (corrected evaluation).}
\label{tab:classification}
\small
\begin{tabular}{@{}llcc@{}}
\toprule
\textbf{Mechanism} & \textbf{Condition} & \textbf{Count} & \textbf{Limit} \\
\midrule
Trivial zero (Thm.~\ref{thm:main}) & $a(1)=0$ & 34 & $K=b(0)$ \\
Finite CF (Prop.~\ref{prop:finite}) & $a(k)=0,\ k\ge 2$ & 9 & $K\in\mathbb{Q}$ \\
Non-trivial rational & --- & 0 & --- \\
\midrule
\textbf{Total} & & \textbf{43/43} & \\
\bottomrule
\end{tabular}
\end{table}

All 34 trivial-zero cases satisfy $K = b(0)$ exactly, verified by
exact rational arithmetic (\texttt{fractions.Fraction}) and by
backward evaluation from depth~100.

The 9 finite CF families terminate at $a(k)=0$ for $k\in\{2,3\}$,
yielding distinct rational values: $-29/13$, $-12/5$, $-7/4$, $-1$,
$2/3$, $8/7$, $3/2$, $28/9$, $4$.
Each value was verified by exact rational arithmetic
(\texttt{fractions.Fraction}) at the termination depth.

No family produced a non-trivial rational limit (i.e., $K\in\mathbb{Q}$
with $K\ne b(0)$ and no finite termination).
No family that was rational under the bugged evaluator became
irrational after correction.

\subsection{Density model}\label{sec:density}

The probability that a random family has $a(1)=0$ is
$\Pr\!\bigl(\sum_{i=0}^d \alpha_i = 0\bigr)$, where each $\alpha_i$ is
uniform on $\{-4,\ldots,4\}$ (the leading coefficient restricted to
$\{-4,\ldots,4\}\setminus\{0\}$).
The leading coefficient $\alpha_d$ is restricted to
$\{-4,\ldots,4\}\setminus\{0\}$, contributing variance $15/2$
rather than the $20/3$ of the remaining coefficients.
By a CLT approximation with total variance $\sigma^2 = d\cdot\tfrac{20}{3}+\tfrac{15}{2}$,
\[
  \Pr(a(1)=0) \;\approx\; \frac{1}{\sqrt{2\pi\sigma^2}},
\]
yielding $\approx 8.7\%$ at $d=2$ and $\approx 5.8\%$ at $d=6$.
Empirically: $7.5\%$ ($d=2$), $6.5\%$ ($d=3$), $8.5\%$ ($d=4$),
$8.5\%$ ($d=5$), $4.5\%$ ($d=6$), consistent within sampling noise.
The finite CF mechanism adds $\sim$1--2\% to the total contamination rate.

% =====================================================================
\section{Pre-screening protocol}\label{sec:protocol}
% =====================================================================

Based on the analysis above, we recommend the following pre-screening
steps before running extended-basis PSLQ on a PCF family:

\begin{enumerate}
  \item \textbf{Evaluate $a(1)$.}
    If $a(1) = \sum\alpha_i = 0$, the family is trivially rational
    ($K=b(0)$).  Discard.

  \item \textbf{Check for integer roots of $a$.}
    If $a(k)=0$ for any $k\ge 2$, the CF terminates at finite depth.
    Compute the finite CF value; if rational, discard or flag as known.

  \item \textbf{Rationality pre-check.}
    Before enriched-basis PSLQ
    (e.g., basis $[1, K, \ln 2, K\ln 2, (\ln 2)^2]$),
    first test the minimal basis $[1, K]$.  If PSLQ returns a relation,
    $K$ is likely rational and will contaminate any enriched scan.

  \item \textbf{Convention audit.}
    Verify the CF evaluator against a known identity
    (e.g., $1+1/(1+1/(1+\cdots)) = \varphi$)
    before trusting numerical values.
\end{enumerate}

Steps~1--2 are $O(d)$ coefficient checks requiring no high-precision
arithmetic.
Step~3 adds a single PSLQ call per family.
Together they eliminated 100\% of false positives in our survey.%
\footnote{The pre-screen is implemented in
  \texttt{scripts/ln2\_prefilter.py}.}

% =====================================================================
\section{Discovery narrative}\label{sec:discovery}
% =====================================================================

The results above were reached through five iterations of an automated
conjecture refinement relay (AEAL architecture~\cite{aeal-paper7,aeal-paper8}),
providing a documented case study of machine-assisted self-correction:

\begin{enumerate}
  \item \textbf{Iteration~7} (false positive):
    A PSLQ scan with the enriched basis
    $[1,\, K,\, \ln 2,\, K\!\cdot\!\ln 2,\, (\ln 2)^2]$
    found 43~hits, initially claimed as evidence that $\ln 2$
    appears in PCF limits at positive density across all degrees.

  \item \textbf{Iteration~8} (rationality diagnosis):
    All 43 relations have the form $[a, b, 0, 0, 0]$
    or factor as $(c_1 K + c_0)(1{+}\ln 2) = 0$,
    both implying $K\in\mathbb{Q}$.
    Zero genuine $\ln 2$ identities.
    The~44 associated claims were retracted.

  \item \textbf{Iteration~9} (partial model):
    Seven families with $K{=}{-}3$ were characterized;
    all share $b(0)+b(1)=-3$ with
    $b(0) \in \{-4,-3,-2,-1\}$.
    A ``dynamical attractor'' hypothesis was proposed
    and later superseded.

  \item \textbf{Iteration~10} (apparent theorem):
    The pattern $a(1) = \sum\alpha_i = 0$ was identified across all
    34 non-finite families.  A Wallis induction proof was constructed
    for the claim ``$K = b(0)+b(1)$.''  This was accepted as a theorem.

  \item \textbf{Iteration~11} (bug discovery):
    A convention audit (Remark~\ref{rem:bug}) revealed that the
    evaluator adds $+b(1)$ to every PCF value.  The ``theorem''
    $K=b(0)+b(1)$ was a documentation of the code bug.
    The corrected result is the trivial $K=b(0)$.
    The Wallis induction proof was valid \emph{within the bugged
    recurrence} but does not describe the standard mathematical
    continued fraction.
\end{enumerate}

The full arc---false positive $\to$ rationality detection $\to$
attractor hypothesis $\to$ apparent theorem $\to$ bug discovery
$\to$ trivial explanation---is, to our knowledge, the first documented
instance of an automated relay generating a false positive,
proposing and then retracting an intermediate ``theorem,''  and
ultimately converging on the correct (trivial) result
through self-correction.
To our knowledge, this is a documented instance of such a complete
arc within an automated relay.
All iterations, including retractions, are preserved in the
governance log.%
\footnote{AEAL claims log: \texttt{results/claims.jsonl}.
  Convention audit: \texttt{scripts/iter11\_bugfix\_reeval.py}.}

% =====================================================================
\section{Discussion}\label{sec:discussion}
% =====================================================================

\textbf{Why this matters.}
The contamination rate of $\sim$5--9\% is high enough to produce
false positives in any PSLQ scan of moderate size.
In our survey, 43 out of 1000 families ($4.3\%$) were affected.
For larger coefficient ranges or higher degrees, the rate changes
predictably via the density model, allowing experimenters to
estimate their expected false-positive load \emph{a priori}.

\textbf{The value of documenting errors.}
The Wallis initialization bug is a cautionary example.
It produces a systematic offset that is invisible to convergence
diagnostics and rationality checks, yet transforms a trivial
observation ($a(1)=0 \Rightarrow K=b(0)$) into an apparently
non-trivial theorem ($K=b(0)+b(1)$) that survives induction
proof---because the proof is valid for the bugged recurrence.
The only reliable diagnostic was cross-checking against a known
identity with an independent evaluator.
We recommend this as standard practice for any computational
continued fraction research.

\textbf{Relation to known results.}
The observation that $a(1)=0$ makes the GCF trivial is implicit
in the definition---the first partial numerator is zero, so the
tail contributes nothing---but does not appear as an explicit
warning in the standard references~\cite{lorentzen-waadeland},
likely because analytic convergence theory assumes $a(n)\ne 0$
for all~$n$.
In the context of random polynomial PCFs, where $a(1)=0$ occurs
at positive density, the observation becomes practically relevant.

\textbf{Broader impact.}
The contamination mechanism identified here---overcomplete PSLQ bases
combined with unvetted evaluator initialization---is not specific to
our relay.
Any system that runs PSLQ on programmatically generated continued
fractions faces the same risk, including the Ramanujan
Machine~\cite{raayoni2021} and its derivatives.
The pre-screening protocol of Section~\ref{sec:protocol} is directly
applicable: steps~1--2 require only coefficient arithmetic and can
be added to any PCF search pipeline at negligible cost.
We hope that documenting both the contamination mechanisms and the
process by which they were discovered---including the intermediate
errors---provides a useful template for quality control in
automated mathematical discovery.

% =====================================================================
% References
% =====================================================================
\begin{thebibliography}{9}

\bibitem{raayoni2021}
G.~Raayoni et al.,
\emph{Generating conjectures on fundamental constants with the
  Ramanujan Machine},
Nature \textbf{590} (2021), 67--73.

\bibitem{lorentzen-waadeland}
L.~Lorentzen and H.~Waadeland,
\emph{Continued Fractions with Applications},
North-Holland, 1992.

\bibitem{aeal-paper7}
Papanokechi,
\emph{AEAL relay architecture for automated conjecture refinement},
Zenodo, 2026.
\texttt{doi:\href{https://doi.org/10.5281/zenodo.19562751}{10.5281/zenodo.19562751}}.

\bibitem{aeal-paper8}
Papanokechi,
\emph{Adversarial epistemic audit loops: governance and self-correction},
Zenodo, 2026.
\texttt{doi:\href{https://doi.org/10.5281/zenodo.19565086}{10.5281/zenodo.19565086}}.

\end{thebibliography}

% =====================================================================
\section*{Declaration on the Use of AI Tools}
% =====================================================================

The computational experiments, verification scripts,
and iterative workflow underlying this paper were
carried out using a multi-agent research environment
combining GitHub Copilot (Microsoft, within Visual
Studio Code) for code generation, script execution,
and autonomous relay iteration, and Claude
(\texttt{claude-sonnet-4-6}, Anthropic) for
high-level mathematical review, proof critique,
and manuscript preparation. The relay architecture
is described in \cite{aeal-paper7,aeal-paper8}.

All mathematical claims, proofs, and numerical
verifications were independently confirmed by the
author. The AI tools were used as research and
writing instruments; all scientific judgements,
including the decision to retract intermediate
claims (Section~\ref{sec:discovery}), were made by the author.

All scripts, data, and the full AEAL governance
log supporting the results are publicly available at:\\
\url{https://github.com/papanokechi/pcf-rational-contamination}.

\end{document}
